\section{An independent uniform NS--PC comparison}\label{sec:comparison}
The linear PC lower bounds were proved by transporting the PC closure
itself. We now establish a different result: every bounded-degree PC
consequence has an NS identity of degree bounded independently of
universe size. The comparison applies to arbitrary polynomials, not
only invariant targets or refutations. Its qualitative consequence is
that NS and PC assign every fixed uniform request the same type.

The NS-span plateau argument already occurs in Kraj\'i\v cek's
uniform comparison~\cite[Theorem~5.5]{Krajicek}. We give an independent
proof of its required chain bound and carry the comparison through for
the actual consequence spaces. A uniform bound on strict
invariant-subspace chains forces the intersections of two successive
NS certificate spaces with the fixed degree-$d$ polynomial ambient
to agree. Their common
space is then closed under every permitted PC operation. The plateau
may depend on $n$, but the bound on its NS certificate degree does not.

At a degree cutoff, $I_e(n)$ is the linear space of polynomials with
degree-at-most-$e$ NS certificates. It is not an ideal closed under
arbitrary multiplication.
A \emph{strict subspace chain} is a sequence in which each inclusion
is proper. Its bound counts proper increases; consecutive equal
spaces form a plateau. The proof uses deletion to preserve those
increases, then uses a plateau to absorb variable multiplication.
No decomposition into simple representations is required.

\begin{lemma}[Uniform chain bound]\label{lem:chains}
For any bounded basis system $\mathcal U$, strict $S_n$-invariant
subspace chains in $K[\mathcal U_n]$ have a length bound independent of $n$.
\end{lemma}
\begin{proof}
Let $r$ bound the basis support. Over an algebraic closure, choose a
fixed $q=p^a>2r$. At $N=n+q$ with $n\ge4r+1$, deletion is exact at the
doubled budget $2r$. Lemma~\ref{lem:detect} shows that each nonzero
subquotient $B/A$ of $K[\mathcal U_N]$ survives. Thus
$A\subsetneq B$ implies
$\mathcal D(A)\subsetneq\mathcal D(B)$: the whole strict chain becomes
a strict chain of the same length in $K[\mathcal U_n]$.

Repeatedly subtract $q$ while the rank exceeds $M=4r+q$. Every retained
rank in such a step exceeds $4r$, so all steps are eligible. The final
chain has at most $\dim K[\mathcal U_M]$ strict inclusions, because a
smaller-rank basis injects into the rank-$M$ basis. The same bound covers
the initial ranks at most $M$. Each step may choose its own generic
coefficients: the detector applies to arbitrary invariant subspaces,
without a naturality assumption on the chain across ranks. Finally,
scalar extension preserves strict inclusions, giving the result over
the original coefficient field.
\end{proof}

Only a finite exactness threshold at each fixed support is needed for
this argument. In characteristic two, the earlier quadratic construction also suffices. The
bound may be large; its role is uniformity in $n$, not a sharp comparison
of the two minimum degrees.

\begin{theorem}[Uniform comparison for all consequences]\label{thm:comparison}
For every uniform axiom template over a field of prime characteristic
$p$, there is a nondecreasing $F:\NN\to\NN$, independent of $n$ and of
the target, with $F(d)\ge d$ and
\begin{equation}\label{eq:whole-comparison}
                              C_d(n)\subseteq I_{F(d)}(n)
                       \qquad(n,d\ge0).
\end{equation}
Consequently, every polynomial with an ordinary-PC derivation of degree
at most $d$ from $\Gamma_n$ has a total-product NS identity of degree at
most $F(d)$ from the same axioms. For the aggregate profiles of any
stable uniform target family, and for the profiles of any one fixed
target orbit,
\begin{equation}\label{eq:target-comparison}
                      h_{\PC}(n)\le h_{\NS}(n)\le F(h_{\PC}(n)),
                      \qquad F(\infty)=\infty.
\end{equation}
Their NS and PC types in Theorem~\ref{thm:classification} therefore agree.
\end{theorem}
\begin{proof}
Fix $d$. Monomials of degree at most $d$ form a permutation basis of
support at most $kd$. Let $L_d$ bound strict invariant-subspace chains in
$P_{n,\le d}$, uniformly in $n$, by Lemma~\ref{lem:chains}.
Consider the invariant spaces
\[
              W_j=I_{d+j}(n)\cap P_{n,\le d},
                    \qquad 0\le j\le L_d+1.
\]
They form an ascending chain. Some $W_j=W_{j+1}$ with $j\le L_d$,
otherwise there would be more than $L_d$ strict inclusions.

This plateau contains every axiom usable within degree $d$ and is closed
under linear combinations. If $f\in W_j$ and $\deg f\le d-1$, then
multiplication of an NS identity for $f$ by a variable gives
$x_vf\in I_{d+j+1}(n)$. Its degree is at most $d$, hence
$x_vf\in W_{j+1}=W_j$. For $d=0$ there is no permitted multiplication
of a nonzero line; the zero line already lies in $W_j$.
Induction on \emph{any} degree-$d$ PC derivation now places every one of
its lines in $W_j$. In particular,
\[
                          C_d(n)\subseteq I_{d+L_d}(n).
\]
This induction retains cancellations: a lower-degree linear combination
can be multiplied because it lies in the plateau, even if its separate
summands have degree $d$.

Choose
\[
                          F(d)=d+\sum_{u=0}^d L_u.
\]
It is nondecreasing, satisfies $F(d)\ge d$, and gives
\eqref{eq:whole-comparison}. No symmetry of a target was used.
Conversely, expansion of NS identities into monomial--axiom products
gives a same-degree PC derivation, so the first inequality in
\eqref{eq:target-comparison} holds. If the aggregate PC degree is finite,
every requested target belongs to $C_d(n)$ at that budget; the inclusion
puts every target in $I_{F(d)}(n)$. Empty requests have degree zero.
Infinite PC degree forces infinite NS degree by the first inequality.
These observations prove the complete extended-order comparison.

For type agreement write $g=h_{\PC}$ and $h=h_{\NS}$.
The inequalities $g\le h\le F(g)$ transfer eventual boundedness in both
directions. If $g$ takes one finite value $d$ infinitely often, then on
those sizes $h$ lies in the finite interval $[d,F(d)]$; some finite value
of $h$ recurs. A recurring finite value of $h$ likewise forces a recurring
finite value of $g$. Thus the recurring sets are empty together.
If recurring values of $g$ are unbounded, the same finite-interval
argument yields unbounded recurring values of $h$. Conversely, choose a
recurring value of $h$ exceeding $F(B)$. On its infinite level set,
$g>B$ and $g\le h$, so some finite value above $B$ recurs for $g$.
Thus infinitude of the recurring sets agrees as well. Boundedness,
emptiness and infinitude determine exactly the four alternatives.
\end{proof}

\begin{corollary}[Common finite peaks]\label{cor:common-peaks}
If a fixed uniform request has common type T4, there is a nondecreasing
$F$ such that, beyond every bound $B$, some size $n\ge B$ has
\[
 h_{\PC}(n)=d\ge B,\qquad d\in\rec(h_{\PC}),\qquad
 d\le h_{\NS}(n)\le F(d)<\infty,
\]
with $n\le(p+2)(\kappa d+\sigma)$. Both degrees satisfy the rational
linear lower bound in \eqref{eq:peaks} at these same sizes.
\end{corollary}
\begin{proof}
Take the finite PC peaks from Theorem~\ref{thm:classification}.
The pointwise lower comparison gives the NS lower bound, and the uniform
upper comparison gives its finiteness at those sizes.
\end{proof}

The comparison does not identify exact degrees, easy-class values or
recurring degree sets. Its potentially large function $F$ cannot be
inverted to obtain a linear PC lower bound from a linear NS lower bound.
That quantitative result came directly from the PC transfer theorem.
